% TOP_PROBLEM: {"problemId": "problem.borel-normal-number-conjecture-for-irrational-algebraic-numbers", "canonicalTitle": "Borel normal number conjecture for irrational algebraic numbers", "releaseRank": 133, "primaryDomain": "number_theory_arithmetic_geometry", "primaryDomainLabel": "Number theory & arithmetic geometry", "displayStatus": "open", "releaseStatus": "open"}
% !TeX program = lualatex
% Definition and sources only; this file is not an LLM solution attempt.
\documentclass[11pt]{article}
\usepackage[margin=25mm]{geometry}
\usepackage{fontspec}
\setmainfont{Latin Modern Roman}
\usepackage{amsmath,amssymb,mathtools}
\usepackage{unicode-math}
\setmathfont{Latin Modern Math}
\usepackage{xurl,hyperref}
\hypersetup{colorlinks=true,urlcolor=blue}
\setcounter{secnumdepth}{0}
\setlength{\parindent}{0pt}
\setlength{\parskip}{5pt}
\setlength{\emergencystretch}{3em}
\begin{document}
\section*{\texorpdfstring{Borel normal number conjecture for irrational algebraic numbers}{Borel normal number conjecture for irrational algebraic numbers}}\label{133}
\subsection{Definitions and mathematical statement}
Write the fractional part of a real number in base $b\ge2$ as $0.d_1d_2\ldots$. It is normal to base $b$ if for every length $k$ and every word $w\in\{0,\ldots,b-1\}^k$,
\[
 \lim_{N\to\infty}\frac{\#\{1\le j\le N:(d_j,\ldots,d_{j+k-1})=w\}}{N}=b^{-k}.
\]
The conjecture requires this for every real algebraic irrational number and every integer base. Algebraic means a root of a nonzero rational polynomial. Frequency of individual digits alone is weaker than normality of all blocks.
\par\smallskip\textit{Formulation and concept references:} \href{https://webusers.imj-prg.fr/~michel.waldschmidt/articles/pdf/WordsTranscendence.pdf}{[S1]}.

\subsection{Short English statement}
Do the digits of every algebraic irrational have the uniform block frequencies of a normal number in every integer base?
\subsection{Sources}
\begin{itemize}
\item[S1]Michel Waldschmidt, Words and Transcendence, 2005.
\url{https://webusers.imj-prg.fr/~michel.waldschmidt/articles/pdf/WordsTranscendence.pdf}.
\textit{Formulation source.}
\end{itemize}
\end{document}
